\documentclass[12pt,a4paper]{article}
\usepackage[utf8]{inputenc}
\usepackage[T5]{fontenc} % Hỗ trợ tiếng Việt
\usepackage{amsmath, amssymb}
\usepackage{mathptmx} % Thiết lập phông chữ Times
\usepackage{geometry}
\geometry{a4paper, margin=2cm}
\usepackage{setspace}
\setstretch{1.3} % Khoảng cách dòng

% --- CẤU HÌNH LOGO CHỮ K IN NHẠT DƯỚI NỀN TRANG ---
\usepackage{background}
\backgroundsetup{
  scale=15,           % Độ to của chữ K (tăng/giảm tùy ý)
  color=gray,         % Tông màu xám
  opacity=0.1,        % Độ nhạt (0.1 là rất mờ và lặn dưới chữ, không lo bị đè)
  angle=0,            % Giữ thẳng đứng không xoay nghiêng
  contents={\textit{K}} % Định dạng chữ K in nghiêng giống mẫu
}

\setlength{\parindent}{0pt} % Xóa thụt lề đầu dòng
\setlength{\parskip}{0.25cm} % Khoảng cách giữa các đoạn văn

\begin{document}

\textbf{Đề bài:} Trong không gian các ma trận vuông cấp hai, xét xem hệ 3 vectơ $A, B, C$ sau đây độc lập tuyến tính hay phụ thuộc tuyến tính:
$$A = \begin{bmatrix} 1 & 2 \\ 0 & 1 \end{bmatrix}, \quad B = \begin{bmatrix} 0 & 2 \\ 1 & 4 \end{bmatrix}, \quad C = \begin{bmatrix} 2 & 6 \\ 1 & 6 \end{bmatrix}$$

Xét phương trình tổ hợp tuyến tính của hệ vectơ bằng vectơ không :
$$x_1A + x_2B + x_3C = \mathbf{0}$$
Thay các ma trận vào, ta được:
$$x_1 \begin{bmatrix} 1 & 2 \\ 0 & 1 \end{bmatrix} + x_2 \begin{bmatrix} 0 & 2 \\ 1 & 4 \end{bmatrix} + x_3 \begin{bmatrix} 2 & 6 \\ 1 & 6 \end{bmatrix} = \begin{bmatrix} 0 & 0 \\ 0 & 0 \end{bmatrix}$$
Cộng các ma trận ở vế trái và đồng nhất các phần tử tương ứng, ta thu được hệ phương trình thuần nhất:
$$\begin{cases}
x_1 + 2x_3 = 0 \quad (1) \\
2x_1 + 2x_2 + 6x_3 = 0 \quad (2) \\
x_2 + x_3 = 0 \quad (3) \\
x_1 + 4x_2 + 6x_3 = 0 \quad (4)
\end{cases}$$
Giải hệ phương trình:\\
Từ $(1) \Rightarrow x_1 = -2x_3$\\
Từ $(3) \Rightarrow x_2 = -x_3$\\
Thay $x_1$ và $x_2$ vào phương trình $(2)$: $2(-2x_3) + 2(-x_3) + 6x_3 = -4x_3 - 2x_3 + 6x_3 = 0$ (Luôn đúng).\\
Thay $x_1$ và $x_2$ vào phương trình $(4)$: $(-2x_3) + 4(-x_3) + 6x_3 = -6x_3 + 6x_3 = 0$ (Luôn đúng).\\
Hệ phương trình có vô số nghiệm phụ thuộc vào $x_3$.\\
Vì hệ phụ thuộc vào $x_3$ nên chỉ cần chọn $x_3$ bất kì thì ta luôn có 1 nghiệm không tầm thường $\Rightarrow$ Tồn tại các hệ số không đồng thời bằng $0$ sao cho tổ hợp tuyến tính bằng $\mathbf{0}$, nên hệ vectơ $\{A, B, C\}$ là phụ thuộc tuyến tính.

\vspace{0.8cm}

\textbf{Giải Bài 3}\\
\textbf{Đề bài:} Cho tập hợp $S = \{(\alpha - \beta, 2\alpha + \beta, \alpha, \beta) \mid \alpha, \beta \in \mathbb{R}\}$. Chứng minh $S$ là không gian con của $\mathbb{R}^4$. Tìm cơ sở và số chiều của $S$.

\textbf{1. Chứng minh $S$ là không gian con:}\\
$\forall u \in S$ đều có dạng $u = (\alpha - \beta, 2\alpha + \beta, \alpha, \beta)$. Ta có thể phân tích:
$$u = (\alpha, 2\alpha, \alpha, 0) + (-\beta, \beta, 0, \beta)$$
$$ \Rightarrow u = \alpha(1, 2, 1, 0) + \beta(-1, 1, 0, 1)$$
Đặt $v_1 = (1, 2, 1, 0)$ và $v_2 = (-1, 1, 0, 1)$. Khi đó, mọi vectơ của $S$ đều là tổ hợp tuyến tính của $v_1$ và $v_2$.\\
Nói cách khác, $S = \text{span}\{v_1, v_2\}$.\\
Vì $S$ là không gian sinh bởi một tập hợp các vectơ trong $\mathbb{R}^4$, nên theo định nghĩa, $S$ là một không gian con của $\mathbb{R}^4$.

\textbf{2. Tìm cơ sở và số chiều của $S$:}\\
Hệ sinh của $S$ là $F = \{v_1, v_2\}$. Xét phương trình:
$$k_1 v_1 + k_2 v_2 = (0,0,0,0)$$
$$\Rightarrow (k_1 - k_2, \quad 2k_1 + k_2, \quad k_1, \quad k_2) = (0,0,0,0)$$
$\Rightarrow k_1 = 0$ và $k_2 = 0$.\\
$\Rightarrow$ hệ $F = \{v_1, v_2\}$ độc lập tuyến tính.\\
Một cơ sở của không gian con $S$ là $B = \{(1, 2, 1, 0), (-1, 1, 0, 1)\}$. Số chiều $\dim(S) = 2$.

\vspace{0.8cm}

\textbf{Giải Bài 11}\\
\textbf{Đề bài:} Không gian $P_3$ gồm đa thức bậc $\le 3$. Xét hai hệ vectơ $(\alpha)$ và $(\beta)$.

\textbf{a) Chứng minh hệ $(\alpha)$ và $(\beta)$ là hai cơ sở của $P_3$:}\\
Không gian $P_3$ có số chiều $\dim(P_3) = 4$.\\
Hệ $(\alpha) = \{1, x, x^2, x^3\}$ gồm $4$ đa thức độc lập tuyến tính (cơ sở chính tắc của $P_3$), do đó $(\alpha)$ là một cơ sở.\\
Để chứng minh $(\beta)$ là cơ sở, ta biểu diễn các đa thức của $(\beta)$ qua cơ sở $(\alpha)$:
$$\beta_1 = 1 = \mathbf{1}\cdot 1$$
$$\beta_2 = x - 2 = \mathbf{-2}\cdot 1 + \mathbf{1}\cdot x$$
$$\beta_3 = (x - 2)^2 = x^2 - 4x + 4 = \mathbf{4}\cdot 1 \mathbf{- 4}\cdot x + \mathbf{1}\cdot x^2$$
$$\beta_4 = (x - 2)^3 = x^3 - 6x^2 + 12x - 8 = \mathbf{-8}\cdot 1 + \mathbf{12}\cdot x \mathbf{- 6}\cdot x^2 + \mathbf{1}\cdot x^3$$
Ma trận tọa độ theo cột của hệ $(\beta)$ đối với cơ sở $(\alpha)$ là:
$$P = \begin{bmatrix} 1 & -2 & 4 & -8 \\ 0 & 1 & -4 & 12 \\ 0 & 0 & 1 & -6 \\ 0 & 0 & 0 & 1 \end{bmatrix}$$
Có định thức $\det(P) = 1 \cdot 1 \cdot 1 \cdot 1 = 1 \neq 0$. Suy ra hệ $(\beta)$ độc lập tuyến tính. Mà hệ có đúng $4$ phần tử bằng số chiều của không gian, nên $(\beta)$ cũng là một cơ sở của $P_3$.

\textbf{b) Tìm ma trận chuyển từ cơ sở $(\alpha)$ sang cơ sở $(\beta)$:}\\
Theo kết quả tính toán ở câu a:
$$P_{\alpha \to \beta} = \begin{bmatrix} 1 & -2 & 4 & -8 \\ 0 & 1 & -4 & 12 \\ 0 & 0 & 1 & -6 \\ 0 & 0 & 0 & 1 \end{bmatrix}$$

\textbf{c) Tìm tọa độ của vectơ $v = x^3 - 3x + 2$ đối với cơ sở $(\beta)$:}\\
Giả sử tọa độ của $v$ đối với $(\beta)$ là $[v]_{\beta} = (c_1, c_2, c_3, c_4)^T$. Khi đó:
$$x^3 - 3x + 2 = c_1(1) + c_2(x-2) + c_3(x-2)^2 + c_4(x-2)^3$$
Khai triển vế phải theo lũy thừa của $x$ và nhóm lại:
$$x^3 - 3x + 2 = c_4 x^3 + (c_3 - 6c_4)x^2 + (c_2 - 4c_3 + 12c_4)x + (c_1 - 2c_2 + 4c_3 - 8c_4)$$
Đồng nhất hệ số hai vế, ta được hệ phương trình:
$$\begin{cases}
c_4 = 1 \\
c_3 - 6c_4 = 0 \\
c_2 - 4c_3 + 12c_4 = -3 \\
c_1 - 2c_2 + 4c_3 - 8c_4 = 2
\end{cases}$$
Giải từ trên xuống dưới:\\
$c_4 = 1$\\
$c_3 = 6(1) = 6$\\
$c_2 = -3 + 4(6) - 12(1) = 9$\\
$c_1 = 2 + 2(9) - 4(6) + 8(1) = 4$\\
Vậy tọa độ cần tìm là:
$$[v]_{\beta} = \begin{bmatrix} 4 \\ 9 \\ 6 \\ 1 \end{bmatrix}$$

\vspace{1.5cm}

\textbf{Giải Bài 5}\\
\textbf{a) Chứng minh $f$ là ánh xạ tuyến tính}\\
$\forall x, y \in \mathbb{R}^4: x=(x_1, x_2, x_3, x_4); y=(y_1, y_2, y_3, y_4); \forall k \in \mathbb{R}$\\
$x+y=(x_1+y_1, x_2+y_2, x_3+y_3, x_4+y_4)$, ta có:
\begin{equation}
\begin{aligned}
f(x+y) &= f(x_1+y_1, x_2+y_2, x_3+y_3, x_4+y_4) \\
&= (2(x_1+y_1)+2(x_2+y_2)-(x_3+y_3), (x_2+y_2)+(x_3+y_3)+(x_4+y_4), (x_3+y_3)+2(x_4+y_4)) \\
&= ((2x_1+2x_2-x_3)+(2y_1+2y_2-y_3), (x_2+x_3+x_4)+(y_2+y_3+y_4), (x_3+2x_4)+(y_3+2y_4)) \\
&= (2x_1+2x_2-x_3, x_2+x_3+x_4, x_3+2x_4) + (2y_1+2y_2-y_3, y_2+y_3+y_4, y_3+2y_4) \\
&= f(x) + f(y)
\end{aligned} \tag{1}
\end{equation}

\begin{equation}
\begin{aligned}
f(kx) &= f(kx_1, kx_2, kx_3, kx_4) \\
&= (2(kx_1)+2(kx_2)-(kx_3), (kx_2)+(kx_3)+(kx_4), (kx_3)+2(kx_4)) \\
&= (k(2x_1+2x_2-x_3), k(x_2+x_3+x_4), k(x_3+2x_4)) \\
&= k(2x_1+2x_2-x_3, x_2+x_3+x_4, x_3+2x_4) \\
&= kf(x)
\end{aligned} \tag{2}
\end{equation}
Từ (1) và (2) ta suy ra $f$ là ánh xạ tuyến tính.

\vspace{0.5cm}

\textbf{b) Tìm ma trận của $f$ trong cặp cơ sở chính tắc.}\\
Ta gọi $\{e_1=(1,0,0,0); e_2=(0,1,0,0); e_3=(0,0,1,0); e_4=(0,0,0,1)\}$ là cơ sở chính tắc trong $\mathbb{R}^4$.\\
$\{e'_1=(1,0,0); e'_2=(0,1,0); e'_3=(0,0,1)\}$ là cơ sở chính tắc trong $\mathbb{R}^3$, như vậy:
\begin{align*}
f(e_1) &= f(1,0,0,0) = (2,0,0) = 2e'_1 + 0e'_2 + 0e'_3 \\
f(e_2) &= f(0,1,0,0) = (2,1,0) = 2e'_1 + 1e'_2 + 0e'_3 \\
f(e_3) &= f(0,0,1,0) = (-1,1,1) = -1e'_1 + 1e'_2 + 1e'_3 \\
f(e_4) &= f(0,0,0,1) = (0,1,2) = 0e'_1 + 1e'_2 + 2e'_3
\end{align*}
Suy ra ma trận $A$ (với các cột là tọa độ của ảnh các vectơ cơ sở) là:
$$A = \begin{bmatrix} 
2 & 2 & -1 & 0 \\ 
0 & 1 & 1 & 1 \\ 
0 & 0 & 1 & 2 
\end{bmatrix}$$

\vspace{0.5cm}

\textbf{c) Tìm $\ker(f), \text{Im}(f)$, cơ sở và số chiều của $\ker(f), \text{Im}(f)$.}\\
\textbf{Tìm $\ker(f)$:}\\
Theo định nghĩa, $x \in \ker(f) \Leftrightarrow f(x) = (0,0,0)$.\\
Ta có hệ phương trình thuần nhất:
$$\begin{cases}
2x_1 + 2x_2 - x_3 = 0 \quad &(1) \\
x_2 + x_3 + x_4 = 0 \quad &(2) \\
x_3 + 2x_4 = 0 \quad &(3)
\end{cases}$$
Giải hệ phương trình trên, từ phương trình (3) suy ra $x_3 = -2x_4$.\\
Thay vào phương trình (2): $x_2 - 2x_4 + x_4 = 0 \Rightarrow x_2 = x_4$.\\
Thay vào phương trình (1): $2x_1 + 2(x_4) - (-2x_4) = 0 \Rightarrow 2x_1 + 4x_4 = 0 \Rightarrow x_1 = -2x_4$.\\
Nghiệm của hệ là $x = (-2x_4, x_4, -2x_4, x_4) = x_4(-2, 1, -2, 1)$ với $\forall x_4 \in \mathbb{R}$.\\
Vậy $\ker(f) = \{(-2b, b, -2b, b) \mid b \in \mathbb{R}\}$.\\
Cơ sở của $\ker(f)$ là $\{(-2, 1, -2, 1)\}$ và số chiều $\dim(\ker(f)) = 1$.

\textbf{Tìm $\text{Im}(f)$:}\\
Không gian ảnh $\text{Im}(f)$ được sinh bởi các cột của ma trận $A$.
$$\text{Im}(f) = \text{Span}\{(2,0,0), (2,1,0), (-1,1,1), (0,1,2)\}$$
Ma trận $A$ đã ở dạng bậc thang với 3 phần tử cơ sở khác không (tại đường chéo chính), do đó hạng của ma trận bằng 3.\\
Từ định lý về số chiều: 
$$\dim(\text{Im}(f)) = \dim(\mathbb{R}^4) - \dim(\ker(f)) = 4 - 1 = 3$$
Vì $\dim(\text{Im}(f)) = 3$ bằng với số chiều của không gian đích $\mathbb{R}^3$, nên tập 3 vectơ cột bất kỳ độc lập tuyến tính sẽ tạo thành một cơ sở. Chẳng hạn, một cơ sở của $\text{Im}(f)$ là $\{(2,0,0), (2,1,0), (-1,1,1)\}$.

\end{document}